\begin{textsample}{POS Dim 1 – 2019 – Score 13.00 – 2019/t000014.txt}  \label{ex:f1_pos_059}
I ' m here to talk to you about something important that may be new to you. \textbf{The} governments of \textbf{the} world are about to conduct an unintentional experiment on our climate.
In 2020, new rules will require ships to lower their sulfur \textbf{emissions} by scrubbing their dirty exhaust or switching to cleaner fuels.
For human health, this is really good, but sulfur particles in \textbf{the} \textbf{emission} of ships also have an effect on clouds.
This is a satellite image of marine clouds off \textbf{the} Pacific West Coast of \textbf{the} United States. \textbf{The} streaks in \textbf{the} clouds are created by \textbf{the} exhaust from ships.
Ships ' \textbf{emissions} include both \textbf{greenhouse} gases, which trap heat over long periods of time, and particulates like sulfates that mix with clouds and temporarily make them brighter.
Brighter clouds reflect more sunlight back to space, cooling \textbf{the} climate.
So in fact, humans are currently running two unintentional experiments on our climate.
In \textbf{the} first one, we ' re increasing \textbf{the} concentration of \textbf{greenhouse} gases and gradually warming \textbf{the} earth system.
This works something like a fever in \textbf{the} human body.
If \textbf{the} fever remains low, its effects are mild, but as \textbf{the} fever rises, damage grows more severe and eventually devastating.
We ' re seeing a little of this now.
In our other experiment, we ' re planning to remove a layer of particles that brighten clouds and shield us from some of this warming. \textbf{The} effect is strongest in ocean clouds like these, and scientists expect \textbf{the} reduction of sulfur \textbf{emissions} from ships next year to produce a measurable increase in global warming.
Bit of a shocker?
In fact, most \textbf{emissions} contain sulfates that brighten clouds : \textbf{coal}, diesel exhaust, forest fires.
Scientists estimate that \textbf{the} total cooling effect from \textbf{emission} particles, which they call aerosols when they ' re in \textbf{the} climate, may be as much as all of \textbf{the} warming we ' ve experienced up until now.
There ' s a lot of uncertainty around this effect, and it ' s one of \textbf{the} major reasons why we have difficulty predicting climate, but this is cooling that we ' ll lose as \textbf{emissions} fall.
So to be clear, humans are currently cooling \textbf{the} \textbf{planet} by dispersing particles into \textbf{the} atmosphere at massive scale.
We just don ' t know how much, and we ' re doing it accidentally.
That ' s worrying, but it could mean that we have a fast-acting way to reduce warming, emergency medicine for our climate fever if we needed it, and it ' s a medicine with origins in nature.
This is a NASA simulation of earth ' s atmosphere, showing clouds and particles moving over \textbf{the} \textbf{planet}. \textbf{The} brightness is \textbf{the} Sun ' s light reflecting from particles in clouds, and this reflective shield is one of \textbf{the} primary ways that nature keeps \textbf{the} \textbf{planet} cool enough for humans and all of \textbf{the} life that we know.
In 2015, scientists assessed possibilities for rapidly cooling climate.
They discounted things like mirrors in space, ping-pong balls in \textbf{the} ocean, plastic sheets on \textbf{the} Arctic, and they found that \textbf{the} most viable approaches involved slightly increasing this atmospheric reflectivity.
In fact, it ' s possible that reflecting just one or two percent more sunlight from \textbf{the} atmosphere could offset two degrees Celsius or more of warming.
Now, I ' m a technology executive, not a scientist.
About a decade ago, concerned about climate, I started to talk with scientists about potential countermeasures to warming.
These conversations grew into collaborations that became \textbf{the} Marine Cloud Brightening Project, which I ' ll talk about momentarily, and \textbf{the} nonprofit policy organization SilverLining, where I am today.
I work with politicians, researchers, members of \textbf{the} tech industry and others to talk about some of these ideas.
Early on, I met British atmospheric scientist John Latham, who proposed cooling \textbf{the} climate \textbf{the} way that \textbf{the} ships do, but with a natural source of particles : sea-salt mist from seawater sprayed from ships into areas of susceptible clouds over \textbf{the} ocean. \textbf{The} approach became known by \textbf{the} name I gave it then,"marine cloud brightening."Early modeling studies suggested that by deploying marine cloud brightening in just 10 to 20 percent of susceptible ocean clouds, it might be possible to offset as much as two degrees Celsius ' s warming.
It might even be possible to brighten clouds in local regions to reduce \textbf{the} impacts caused by warming ocean \textbf{surface} temperatures.
For example, regions such as \textbf{the} Gulf Atlantic might be cooled in \textbf{the} months before a hurricane season to reduce \textbf{the} force of storms.
Or, it might be possible to cool waters flowing onto coral reefs overwhelmed by heat stress, like Australia ' s Great Barrier Reef.
But these ideas are only theoretical, and brightening marine clouds is not \textbf{the} only way to increase \textbf{the} reflection of \textbf{the} sunlight from \textbf{the} atmosphere.
Another \textbf{occurs} when large volcanoes release material with enough force to reach \textbf{the} upper layer of \textbf{the} atmosphere, \textbf{the} stratosphere.
When Mount Pinatubo erupted in 1991, it released material into \textbf{the} stratosphere, including sulfates that mix with \textbf{the} atmosphere to reflect sunlight.
This material remained and circulated around \textbf{the} \textbf{planet}.
It was enough to cool \textbf{the} climate by over half a degree Celsius for about two years.
This cooling led to a striking increase in Arctic ice cover in 1992, which dropped in subsequent years as \textbf{the} particles fell back to earth.
But \textbf{the} volcanic phenomenon led Nobel Prize winner Paul Crutzen to propose \textbf{the} idea that dispersing particles into \textbf{the} stratosphere in a controlled way might be a way to counter global warming.
Now, this has risks that we don ' t understand, including things like heating up \textbf{the} stratosphere or damage to \textbf{the} ozone layer.
Scientists think that there could be safe approaches to this, but is this really where we are?
Is this really worth considering?
This is a simulation from \textbf{the} US National Center for Atmospheric Research global climate model showing, earth \textbf{surface} temperatures through 2100. \textbf{The} globe on \textbf{the} \textbf{left} visualizes our current trajectory, and on \textbf{the} right, a world where particles are introduced into \textbf{the} stratosphere gradually in 2020, and maintained through 2100.
Intervention keeps \textbf{surface} temperatures near those of today, while without it, temperatures rise well over three degrees.
This could be \textbf{the} difference between a safe and an unsafe world.
So, if there ' s even a chance that this could be close to reality, is this something we should consider seriously?
Today, there are no capabilities, and scientific knowledge is extremely limited.
We don ' t know whether these types of interventions are even feasible, or how to characterize their risks.
Researchers hope to explore some basic questions that might help us know whether or not these might be real options or whether we should rule them out.
It requires multiple ways of studying \textbf{the} climate system, including \textbf{computer} models to forecast changes, analytic \textbf{techniques} like machine \textbf{learning}, and many types of observations.
And though it ' s controversial, it ' s also critical that researchers develop core technologies and perform small-scale, real-world experiments.
There are two research programs proposing experiments like this.
At Harvard, \textbf{the} SCoPEx experiment would release very small amounts of sulfates, calcium carbonate and water into \textbf{the} stratosphere with a balloon, to study \textbf{chemistry} and physics effects.
How much material?
Less than \textbf{the} amount released in one minute of flight from a commercial aircraft.
So this is definitely not dangerous, and it may not even be scary.
At \textbf{the} University of Washington, scientists hope to spray a fine mist of salt water into clouds in a series of land and ocean tests.
If those are successful, this would culminate in experiments to measurably brighten an area of clouds over \textbf{the} ocean. \textbf{The} marine cloud brightening effort is \textbf{the} first to develop any technology for generating aerosols for atmospheric sunlight reflection in this way.
It requires producing very tiny particles — think about \textbf{the} mist that comes out of an asthma inhaler — at massive scale — so think of looking up at a cloud.
It ' s a tricky engineering problem.
So this one nozzle they developed generates three trillion particles per second, 80 nanometers in size, from very corrosive saltwater.
It was developed by a team of retired engineers in Silicon Valley — here they are — working full-time for six years, without pay, for their grandchildren.
It will take a few million dollars and another year or two to develop \textbf{the} full spray system they need to do these experiments.
In other parts of \textbf{the} world, research efforts are emerging, including small modeling programs at Beijing Normal University in China, \textbf{the} Indian Institute of Science, a proposed center for climate repair at Cambridge University in \textbf{the} UK and \textbf{the} DECIMALS Fund, which sponsors researchers in global South countries to study \textbf{the} potential impacts of these sunlight interventions in their part of \textbf{the} world.
But all of these programs, including \textbf{the} experimental ones, lack significant funding.
And understanding these interventions is a hard problem. \textbf{The} earth is a vast, complex system and we need major investments in climate models, observations and basic science to be able to predict climate much better than we can today and manage both our accidental and any intentional interventions.
And it could be urgent.
Recent scientific reports predict that in \textbf{the} next few decades, earth ' s fever is on a path to devastation : extreme heat and fires, major loss of ocean life, collapse of Arctic ice, displacement and suffering for hundreds of millions of people. \textbf{The} fever could even reach tipping points where warming takes over and human efforts are no longer enough to counter accelerating changes in natural systems.
To prevent this circumstance, \textbf{the} UN ' s International Panel on Climate Change predicts that we need to stop and even reverse \textbf{emissions} by 2050.
How?
We have to quickly and radically transform major economic sectors, including energy, construction, \textbf{agriculture}, \textbf{transportation} and others.
And it is imperative that we do this as fast as we can.
But our fever is now so high that climate experts say we also have to remove massive quantities of CO2 from \textbf{the} atmosphere, possibly 10 times all of \textbf{the} world ' s annual \textbf{emissions}, in ways that aren ' t proven yet.
Right now, we have slow-moving solutions to a fast-moving problem.
Even with \textbf{the} most optimistic assumptions, our exposure to risk in \textbf{the} next 10 to 30 years is unacceptably high, in my opinion.
Could interventions like these provide fast-acting medicine if we need it to reduce \textbf{the} earth ' s fever while we address its underlying causes?
There are real concerns about this idea.
Some people are very worried that even researching these interventions could provide an excuse to delay efforts to reduce \textbf{emissions}.
This is also known as a moral hazard.
But, like most medicines, interventions are more dangerous \textbf{the} more that you do, so research actually tends to draw out \textbf{the} fact that we absolutely, positively cannot continue to fill up \textbf{the} atmosphere with \textbf{greenhouse} gases, that these kinds of alternatives are risky and if we were to use them, we would need to use as little as possible.
But even so, could we ever learn enough about these interventions to manage \textbf{the} risk?
Who would make decisions about when and how to intervene?
What if some people are worse off, or they just think they are?
These are really hard problems.
But what really worries me is that as climate impacts worsen, leaders will be called on to respond by any means available.
I for one don ' t want them to act without real information and much better options.
Scientists think it will take a decade of research just to assess these interventions, before we ever were to develop or use them.
Yet today, \textbf{the} global level of investment in these interventions is effectively zero.
So, we need to move quickly if we want policymakers to have real information on this kind of emergency medicine.
There is hope! \textbf{The} world has solved these kinds of problems before.
In \textbf{the} 1970s, we identified an existential threat to our protective ozone layer.
In \textbf{the} 1980s, scientists, politicians and industry came together in a solution to replace \textbf{the} chemicals causing \textbf{the} problem.
They achieved this with \textbf{the} only legally binding environmental agreement signed by all countries in \textbf{the} world, \textbf{the} Montreal Protocol.
Still in force today, it has resulted in a recovery of \textbf{the} ozone layer and is \textbf{the} most successful environmental protection effort in human history.
We have a far greater threat now, but we do have \textbf{the} ability to develop and agree on solutions to protect people and restore our climate to health.
This could mean that to remain safe, we reflect sunlight for a few decades, while we green our industries and remove CO2.
It definitely means we must work now to understand our options for this kind of emergency medicine.
Thank you,

% matched lemmas: agriculture, chemistry, coal, computer, emission, greenhouse, learning, left, occur, planet, surface, technique, the, transportation
\end{textsample}
